% native TikZ fragment (batch_095, kind=single_function_options)
% 2013 山东文 9 (A)  y=x cos x + sin x
%% shape A : odd; f(0)=0, f'(0)~1.44>0; single max ~1.55 at x~1.70 (min mirrored); no further zero; f(pi)~0.95>0; f(3.6)~0.92>0
\usetikzlibrary{arrows.meta}
\begin{tikzpicture}[x=1cm,y=1cm,line width=.6pt,>=Stealth,font=\normalsize]
\draw[->] (-3.95,0)--(3.95,0) node[right,inner sep=0pt,xshift=2pt]{$x$};
\draw[->] (0,-3.95)--(0,3.95) node[above,inner sep=0pt,yshift=2pt]{$y$};
\draw (3.1416,-0.10)--(3.1416,0.10);
\node[anchor=north,inner sep=0pt] at (3.1416,-0.26) {$\pi$};
\node[anchor=north west,inner sep=0pt,xshift=2pt,yshift=-2pt] at (0,0) {$O$};
\draw[domain=-3.60:3.60,samples=281] plot (\x,{1.7901*sin(deg(\x))-0.3073*\x+0.0611*\x*\x*\x});
\end{tikzpicture}
